\documentclass[11pt,a4paper]{article}
\usepackage[margin=1.8cm]{geometry}
\usepackage[T1]{fontenc}
\usepackage[utf8]{inputenc}
\usepackage{lmodern}
\usepackage{tabularx}
\usepackage{array}
\usepackage[table]{xcolor}
\usepackage{enumitem}
\usepackage[normalem]{ulem}

\definecolor{sgreen}{HTML}{2E7D32}
\definecolor{wred}{HTML}{C62828}
\definecolor{oblue}{HTML}{1565C0}
\definecolor{torange}{HTML}{E65100}

\setlist[itemize]{leftmargin=1.2em, itemsep=3pt, topsep=2pt, parsep=0pt}
\renewcommand{\arraystretch}{1.4}
\setlength{\parindent}{0pt}
\pagestyle{empty}

\newcommand{\head}[2]{\cellcolor{#1}\textcolor{white}{\textbf{\large #2}}}

\begin{document}

\begin{center}
  {\LARGE\bfseries CD PROJEKT RED -- SWOT Analysis}\\[4pt]
  {\small Polish AAA game studio (Warsaw) \textbullet{} \textit{The Witcher}, \textit{Cyberpunk 2077}, \textit{Gwent}}
\end{center}

\vspace{6pt}

\begin{tabularx}{\textwidth}{|>{\raggedright\arraybackslash}X|>{\raggedright\arraybackslash}X|}
\hline
\head{sgreen}{Strengths} & \head{wred}{Weaknesses} \\
\hline
\begin{itemize}
  \item Two globally known brands: \textit{The Witcher} and \textit{Cyberpunk}
  \item Reputation for strong stories, quests and characters
  \item Large expansions (\textit{Blood and Wine}, \textit{Phantom Liberty}) and long free support
  \item Own store (GOG) and solid finances from older games
  \item Free \textit{Witcher 3} remaster (2022) brought back more players on Steam than at launch
\end{itemize}
&
\begin{itemize}
  \item Damaged trust after the broken \textit{Cyberpunk 2077} launch (2020)
  \item Very long gaps between games; income depends on 1--2 titles
  \item History of forced overtime and overpromising
  \item Loss of own-engine know-how (REDengine) during the switch to Unreal Engine 5
\end{itemize}
\\
\hline
\head{oblue}{Opportunities} & \head{torange}{Threats} \\
\hline
\begin{itemize}
  \item Netflix series and \textit{Edgerunners} anime bring new fans
  \item Steady demand for big single-player story RPGs
  \item New consoles (PS5, Xbox Series, Switch 2) let old games sell again
  \item Unreal Engine 5 makes hiring easier and allows several projects at once
  \item \uline{Booming physical card games: a \textit{Gwent} TCG with booster packs}
  \item \uline{Episodic \textit{Gwent} story campaigns (like \textit{Thronebreaker}) set in the world of \textit{The Witcher 4}; the first episode advertises the main game}
\end{itemize}
&
\begin{itemize}
  \item Players and investors punish broken launches hard
  \item Rising game budgets and strong RPG rivals (Bethesda, BioWare, Larian, Warhorse)
  \item Competition for experienced developers
  \item Delays or tech problems with the new engine and parallel projects
  \item \uline{Small audience for \textit{Gwent} spin-offs: neither \textit{Thronebreaker} nor \textit{Rogue Mage} got a sequel}
  \item \uline{Strong card game rivals (\textit{Hearthstone}, \textit{Balatro}, \textit{Slay the Spire}) and big TCGs (\textit{Pokémon}, \textit{Magic})}
  \item \uline{A delay of \textit{The Witcher 4} also delays the campaigns tied to it}
  \item \uline{Players may see paid episodes and bonus cards as a cash grab}
\end{itemize}
\\
\hline
\end{tabularx}

\end{document}
